\textbf{结论名称：}抛物线切线方程（切点式）\\
\textbf{适用条件：}抛物线为标准方程 \(y^2 = 2px\)（\(p>0\)，开口向右），且点 \(P(x_0, y_0)\) 在抛物线上，即满足 \(y_0^2 = 2p x_0\)。对于开口向左、向上、向下的情形可类比变形。\\
\textbf{变量说明：}\(p\) 为焦准距（\(p>0\)），\((x_0, y_0)\) 为切点坐标。\\
\textbf{核心公式：}\[
\boxed{y_0 y = p(x + x_0)}
\]\\
\textbf{使用场景：}已知抛物线上一点直接写切线方程；求解切点弦问题；涉及抛物线切线性质的定点定值问题。\\
\textbf{简单例子：}已知抛物线 \(y^2 = 4x\) 上一点 \((1,2)\)，求该点处的切线方程。\\
解：由 \(y^2 = 4x\) 得 \(2p = 4\)，故 \(p = 2\)，切点 \((1,2)\)。代入公式：
\[2 \cdot y = 2 \cdot (x + 1) \;\Longrightarrow\; y = x + 1.\]
验证：联立 \(\begin{cases} y^2 = 4x \\ y = x+1 \end{cases}\)，消去 \(y\) 得 \((x+1)^2 = 4x \Rightarrow x^2 - 2x + 1 = 0\)，判别式 \(\Delta = 0\)，直线与抛物线相切，符合要求。